
Code generation models are evaluated on multiple benchmarks—HumanEval, MBPP, APPS—to assess performance. Benchmark papers describe different evaluation philosophies: HumanEval emphasizes "algorithmic clarity" (Chen et al., 2021), MBPP targets "practical programming patterns" (Austin et al., 2021), and APPS focuses on "competitive programming competence" (Hendrycks et al., 2021). These design differences lead to the expectation that benchmarks measure distinct competencies.

No prior work has empirically tested whether execution-based code generation benchmarks produce independent model rankings. Current practice reports scores on multiple benchmarks independently, implicitly treating each as measuring separate evaluation dimensions. This assumption has not been validated through quantitative correlation analysis.

We investigate whether HumanEval and MBPP, despite documented design philosophy differences, produce distinctive evaluation signatures measured by model ranking divergence and distributional differences. Our analysis examines 8 code generation models evaluated on both benchmarks.

The investigation proceeds through two phases: (1) establishing data infrastructure by extracting execution trace features with ≥95\% completeness, and (2) testing ranking distinctiveness through correlation and divergence analysis. The first phase succeeded with 100\% feature completeness. The second phase revealed perfect ranking correlation (ρ = 1.000) despite high distributional divergence (KL = 18.4), indicating identical competency orderings with different statistical properties.

This finding has implications for benchmark selection and evaluation design. If benchmarks rank models identically, using multiple execution-based benchmarks provides redundant ranking information. Evaluation diversity may require including different task types rather than multiple generation benchmarks.

We contribute: (1) validated methodology for extracting standardized execution trace features across benchmarks, (2) empirical evidence of perfect ranking correlation between HumanEval and MBPP, and (3) analysis separating distributional divergence from ranking divergence. Limitations include small sample size (n=6-8 models) and restriction to two benchmarks.
